\chapter{Background and Literature Review}
\label{ch:background}

This chapter surveys three bodies of literature whose convergence establishes
the research gap this dissertation addresses.  The central question is: under
what conditions is reference-label-based evaluation a valid methodology for
comparing automated AI ethics assessors?  Stream~1 examines how existing AI
ethics assessment tools are evaluated and reveals a systematic failure to
check reference-label reliability.  Stream~2 surveys the human-label-variation
literature, which demonstrates that annotator disagreement in value-laden
domains is frequently a signal of genuine ambiguity rather than annotation
error.  Stream~3 examines contestability theory and validity frameworks,
which together supply both the normative justification and the
measurement-theoretic foundation for alternative evaluation criteria when
reference-based evaluation is uninformative.  Each stream closes on the
specific deficiency the diagnostic methodology proposed in this dissertation
is designed to address.

% ------------------------------------------------------------------
\section{Stream 1: Evaluation Methodology and Benchmark Validity}
% ------------------------------------------------------------------

\subsection{How AI Ethics Assessment Tools Are Currently Evaluated}

The small but growing literature on automated ethics and compliance
assessment reveals a consistent evaluation pattern: system outputs are
scored against reference labels provided by the system's developers or a
small panel of experts, using standard classification metrics.

RiskRAG~\cite{riskrag2025}, presented at CHI~2025, evaluates its
risk-assessment outputs against human expert judgements on model cards,
reporting precision and recall of identified risks.  PASTA~\cite{pasta2026},
at CHI~2026, validates system-generated compliance scores against a panel of
human experts' median judgements across multiple policy frameworks --- an
approach directly analogous to the evaluation design used in early versions of
the system studied in this dissertation.  Farber~\cite{farber2025} evaluates
AI-based offender risk assessment against human expert judgements, finding
systematic divergence in which the AI assigns higher risk scores, but frames
this as a calibration finding rather than questioning whether the human
judgements themselves are sufficiently reliable to serve as a reference
standard.

In all three cases, the evaluation assumes that expert labels constitute a
valid reference against which system performance can be meaningfully measured.
None reports the inter-annotator reliability of the reference labels, and none
considers the possibility that low system-versus-expert agreement might
reflect label unreliability rather than system inadequacy.  This omission is
not unique to these papers; it characterises the field.  The consequence is
that two assessors achieving different accuracy scores against the same
reference set may genuinely differ in quality, or the reference set may
simply be too unreliable to discriminate between them.  Without checking,
neither conclusion can be drawn.

\subsection{Benchmark Validity in NLP and Machine Learning}

The assumption of reliable reference labels has been challenged more broadly
in NLP evaluation.  Bowman and Dahl~\cite{bowman2021} argue that evaluation
for natural language understanding is ``broken,'' identifying four criteria
benchmarks should meet --- including annotation reliability --- and showing
that most current benchmarks fail on at least one.  They note that
``adversarial data collection does not meaningfully address the causes of
these failures'' and that restoring a healthy evaluation ecosystem requires
progress in annotation reliability, benchmark size, and handling of social
bias.

Raji, Bender, Paullada, Denton, and Hanna~\cite{raji2021} go further,
questioning the framing of benchmarks as general measures of progress.
Their NeurIPS~2021 paper demonstrates how the AI research community treats
performance on a small collection of influential benchmarks as indicative of
progress toward flexible, generalizable AI --- a framing that obscures the
limited and constructed nature of the evaluation tasks and the assumptions
embedded in their reference labels.

Most directly relevant to this dissertation is the Evidence-Centred Benchmark
Design (ECBD) framework~\cite{ecbd2024}, published at ACL~2024.  ECBD draws
explicitly on evidence-centred design from educational measurement theory ---
the same intellectual tradition as Messick's validity framework
\cite{messick1995} --- and applies it to NLP benchmarks.  Case studies of
BoolQ, SuperGLUE, and HELM reveal ``frequent threats to validity when
ECD principles are not observed, including ill-defined constructs, unjustified
data selection, and opaque aggregation.''  This finding directly supports the
central claim of this dissertation: evaluation protocols are measurement
instruments whose validity conditions should be established before they are
used to compare systems.

\subsection{The Reproducibility of Human Evaluation}

A parallel challenge to the reliability of reference labels comes from the
reproducibility literature.  Belz and colleagues, working within the ADAPT
Centre and Dublin City University, have systematically investigated the
reproducibility of human evaluations of NLP systems.  The ReproHum project
and associated ReproNLP shared tasks~\cite{belz2025gem} demonstrate that
human evaluation results in natural language generation are frequently not
reproducible: different experimental set-ups, annotator pools, and
operationalisations of the same quality dimension yield significantly different
results.

Belz's QCET Taxonomy~\cite{belz2025qcet} shows that NLP evaluation
experiments reporting results for the same quality criterion name --- for
example, \emph{Fluency} --- do not necessarily evaluate the same aspect of
quality, and ``the comparability implied by the name can be misleading.''
The QRA++ framework~\cite{belz2025qra} proposes quantified reproducibility
assessment for common result types in NLP, providing the first systematic
methodology for determining whether a given human evaluation result is likely
to replicate under independent conditions.

Taken together, this body of work establishes that even when human evaluations
exist, they may not be reproducible or commensurable across studies.  If the
same evaluation criterion yields different results under different conditions,
any single set of reference labels represents one of several possible
measurement outcomes --- not the ground truth.  This compounds the
reference-label reliability problem identified above: the issue is not merely
whether labels were collected carefully, but whether the evaluation construct
is stable enough to support the comparisons being drawn.

\subsection{Gap}

Current evaluation practice for AI ethics assessment tools treats expert
labels as ground truth without checking whether those labels are sufficiently
reliable to support the discriminations being drawn between systems.  The
benchmark-validity and reproducibility literatures demonstrate that this
assumption fails across NLP more broadly, but no existing work applies this
insight specifically to AI ethics assessment, where the inherently value-laden
nature of the task makes label unreliability especially likely.  The
diagnostic methodology proposed in this dissertation is designed to fill this
gap by making the reliability check explicit and computable.

% ------------------------------------------------------------------
\section{Stream 2: Human Label Variation --- Disagreement as Signal}
% ------------------------------------------------------------------

\subsection{The Perspectivist Turn in Annotation Research}

A fundamental challenge to reference-label-based evaluation emerges from the
growing recognition that annotator disagreement is not always noise.  Aroyo
and Welty~\cite{aroyo2015} articulate this most directly in ``Truth Is a Lie:
Crowd Truth and the Seven Myths of Human Annotation,'' arguing that human
annotation of semantic interpretation tasks is based on ``an antiquated ideal
of a single correct truth.''  They expose seven myths --- including that there
is one correct interpretation per item, that disagreement indicates poor
annotation quality, and that one label per item suffices --- and propose
CrowdTruth, a framework that treats the distribution of disagreements as
informative data rather than noise to be eliminated.

Plank~\cite{plank2022} provides a comprehensive survey of human label
variation in ``The `Problem' of Human Label Variation: On Ground Truth in
Data, Modeling and Evaluation'' (EMNLP~2022).  The paper reconciles
previously proposed notions of label variation, catalogues datasets retaining
un-aggregated labels, and argues that ``human variation in labelling is often
considered noise, and annotation projects for machine learning aim at
minimising human label variation with the assumption to maximise data quality
and machine learning metrics.  However, this conventional practice assumes a
ground truth exists and neglects genuine human variation.''

Uma, Fornaciari, Hovy, Paun, Plank, and Poesio~\cite{uma2021} survey the
literature on learning from disagreement in the \emph{Journal of Artificial
Intelligence Research}, reviewing evidence across NLP and computer vision tasks
and discussing approaches to training models from datasets containing multiple
judgements in disagreement.  The survey establishes that disagreement is
pervasive even in supposedly objective tasks such as part-of-speech tagging,
and that existing techniques for aggregating annotations often discard
information that would otherwise be available for learning.

\subsection{Inherent Disagreement versus Annotation Error}

A critical distinction for this dissertation is between disagreement that
reflects genuine ambiguity or value pluralism --- irreducible, or aleatoric,
disagreement --- and disagreement that reflects poor task design, unclear
instructions, or annotator incompetence --- reducible, or epistemic,
disagreement.  The two require different responses: reducible disagreement
calls for better annotation procedures; irreducible disagreement calls for
rethinking whether a single correct label exists.

Pavlick and Kwiatkowski~\cite{pavlick2019} demonstrate in ``Inherent
Disagreements in Human Textual Inferences'' (TACL~2019) that disagreements in
natural language inference ``are not dismissible as annotation noise, but
rather persist as more ratings are collected and as the amount of context
provided to raters varies.''  They argue for evaluation that requires models
to ``explicitly capture the full distribution of plausible human judgements''
rather than a single gold label, and provide empirical evidence that the
uncertainty captured by state-of-the-art NLI models does not reflect the type
of uncertainty present in human disagreements.

Basile, Fell, Fornaciari, Hovy, Paun, Plank, Poesio, and Uma~\cite{basile2021}
extend this argument explicitly to evaluation methodology in ``We Need to
Consider Disagreement in Evaluation'' (2021), arguing that ``current
evaluation practice largely hinges on the existence of a single `ground truth'
against which we can meaningfully compare the prediction of a model'' and that
``current methods of adjudication, agreement, and evaluation need serious
reconsideration.''  The information-sufficiency analysis in this dissertation
--- distinguishing proposals where annotators had insufficient information
(epistemic) from proposals where disagreement persisted despite adequate
information (aleatoric) --- operationalises this distinction empirically.

\subsection{Perspective-Dependent Labelling}

Gordon, Lam, Park, Patel, Hancock, Hashimoto, and Bernstein~\cite{gordon2022}
introduce ``Jury Learning'' (CHI~2022), demonstrating that annotation
disagreements in subjective tasks reflect legitimate differences in perspective
rather than errors.  They propose an architecture that models every annotator
individually, allowing practitioners to compose ``juries'' whose composition
determines the classifier's predictions.  This work establishes a precedent
for treating annotator identity --- including professional background and
values --- as a structured variable rather than a source of noise to be
controlled away.

The relevance to AI ethics assessment is direct: a human rights lawyer, a
machine learning engineer, and an affected community member may legitimately
disagree about whether a research proposal poses a ``High'' ethical risk, not
because any of them is wrong, but because they are applying different
evaluative frameworks grounded in different professional norms and values.
The stratified inter-annotator analysis in this dissertation --- computing
within-perspective agreement separately from cross-perspective agreement ---
is designed to detect and quantify exactly this structure.

\subsection{Gap}

While the perspectivist literature has established that disagreement should be
treated as signal rather than noise, and has developed computational methods
for learning from disagreement, it has not addressed the upstream implication
for evaluation methodology: if reference labels are perspective-dependent,
then evaluation metrics computed against any single perspective's labels have
limited validity for all perspectives, and the question of when to trust such
metrics remains open.  This dissertation addresses that question by making
the perspective-structure of disagreement a diagnostic input.

% ------------------------------------------------------------------
\section{Stream 3: Contestability and Evaluation Beyond Accuracy}
% ------------------------------------------------------------------

\subsection{Contestable AI}

When reference-label-based evaluation is uninformative, alternative evaluation
criteria are needed.  The contestable AI literature provides a normatively
grounded candidate.

Alfrink, Keller, Kortuem, and Doorn~\cite{alfrink2023contestable} propose a
design framework in ``Contestable AI by Design'' (\emph{Minds and Machines},
2023), arguing that AI systems affecting people's lives should be designed to
allow those affected to understand, challenge, and seek revision of automated
decisions.  Their framework identifies design principles for contestability
--- transparency of the decision basis, intelligibility of the explanation,
and availability of a revision mechanism --- but does not operationalise these
as measurable evaluation criteria with defined computation procedures.

Henin and Le~M\'{e}tayer~\cite{henin2022} argue in ``Beyond Explainability:
Justifiability and Contestability of Algorithmic Decision Systems''
(\emph{AI~\& Society}, 2022) that explainability is necessary but not
sufficient for the legitimacy of algorithmic decision systems.  They
distinguish two key requirements: \emph{justifiability} --- decisions can be
defended by reference to applicable norms --- and \emph{contestability} ---
affected parties can challenge decisions effectively.  Crucially, they show
that a system can be highly explainable yet still fail to be contestable, if
its explanations do not map onto the normative concepts by which the affected
party would dispute the decision.

Hildebrandt~\cite{hildebrandt2019agonistic} grounds contestability in a
deeper philosophical argument, observing that machine learning systems
construct ``computable selves'' from data that may not correspond to subjects'
own self-understanding, and that the right to contest is a necessary safeguard
against this mismatch between a system's model of a person and that person's
lived experience.

Selbst and Barocas~\cite{selbst2018} distinguish between \emph{inscrutability}
--- inability to understand a model's rules --- and \emph{non-intuitiveness}
--- inability to understand why the rules are what they are --- arguing that
explanations must engage the affected party's intuitions to support genuine
contestability.  This distinction informs the definition of
\emph{comprehensibility} as an evaluation property in Chapter~\ref{ch:project}:
it is not sufficient for an assessor's output to be traceable; it must be
understandable to the person whose proposal is being assessed.

\subsection{Legal and Regulatory Requirements for Contestability}

The normative case for contestability is reinforced by legal requirements that
create a practical demand for evaluation approaches beyond accuracy.  GDPR
Article~22 establishes the right not to be subject to decisions based solely
on automated processing that produce legal or similarly significant effects,
including the right to obtain human intervention and to contest such decisions.
The EU AI Act~\cite{euaiact2024} extends these requirements to high-risk AI
systems, mandating human oversight (Article~14), transparency (Article~13),
and mechanisms through which affected persons may seek effective remedy.

Lasek-Markey and Hogan~\cite{lasekmarkey2025} analyse the adequacy of the
FRIA model under Article~27 of the AI Act, arguing that the Act's enumeration
of 23 rights is too narrow --- all substantive EU Charter rights are
potentially engaged by AI systems --- and identifying limitations in the
FRIA's capacity to identify and mitigate fundamental rights impacts.  Janssen,
Lee, and Singh~\cite{janssen2022} propose a practical four-phased framework
for fundamental rights impact assessment under GDPR, providing operational
guidance for deployers but not addressing how the quality of those assessments
should be measured and compared.

These legal instruments create a practical demand for evaluation approaches
that assess whether automated systems satisfy procedural requirements ---
transparency, contestability, human oversight --- rather than merely whether
their outputs match a reference standard.  An assessor that satisfies these
procedural requirements has legitimate governance value even in the absence of
a reliable reference against which its accuracy can be computed.

\subsection{Validity Theory as Theoretical Foundation}

The connection between reference-label unreliability and the need for
alternative evaluation criteria is supplied by validity theory from educational
measurement.

Messick~\cite{messick1995} redefines validity as a unified concept with six
aspects, of which the most relevant here is \emph{consequential validity}: the
appropriateness and social consequences of score interpretation and use.  In
Messick's framework, when no stable criterion exists against which to validate
a measure, validity cannot rest on criterion validity; it must instead rest on
the justifiability of the measure's consequences.  Applied to AI ethics
assessment: when reference labels are unreliable --- as the diagnostic
methodology in this dissertation is designed to detect --- accuracy-based
evaluation lacks criterion validity.  Alternative evaluation criteria must
then be justified on consequential grounds, asking not ``does the assessor's
output match the label?'' but ``does the assessor's output support decisions
whose consequences are legitimate?''  Contestability is the operationalisation
of that consequential criterion.

The ECBD framework~\cite{ecbd2024} demonstrates that it is feasible and
productive to apply Messick's validity theory to AI evaluation practice.  This
dissertation extends that application from NLP benchmarks to the specific
domain of AI ethics assessment, where the value-laden nature of the task makes
the validity question especially pressing.

\subsection{From Accuracy to Legitimacy in Practice}

Huang~\cite{huang2025moderation} provides the closest practical parallel to
this dissertation's argument in ``Content Moderation by LLM: From Accuracy to
Legitimacy'' (\emph{Artificial Intelligence Review}, 2025).  Huang argues that
accuracy is ``insufficient and misleading'' for evaluating content moderation
systems because it ``fails to grasp the broader requirements for legitimate
moderation systems,'' and proposes a shift toward evaluation frameworks
emphasising procedural legitimacy, transparency, and user participation.
Content moderation and AI ethics assessment share a critical structural feature:
both involve value-laden classification where ``correct'' is genuinely
contested, and both deploy AI systems whose decisions affect the rights and
opportunities of individuals who have no meaningful recourse if the only
evaluation criterion is aggregate accuracy against opaque reference labels.

\subsection{Gap}

The contestability literature provides normatively grounded evaluation
criteria, and validity theory provides the formal justification for when those
criteria should replace accuracy-based evaluation.  However, no existing work
combines these into a diagnostic procedure that first determines whether
accuracy-based evaluation is valid for a specific assessor--corpus pair, and
then prescribes alternative evaluation when it is not.  That combination is
the contribution of this dissertation.

% ------------------------------------------------------------------
\section{Summary and Research Gap}
% ------------------------------------------------------------------

Three converging bodies of literature establish the conditions for the
contribution of this dissertation.

\textbf{Evaluation methodology} reveals that AI ethics assessment tools are
evaluated against reference labels whose reliability is never checked, and
that the broader NLP community has identified benchmark validity as a systemic
problem requiring measurement-theoretic solutions~\cite{bowman2021, raji2021,
ecbd2024, belz2025qcet}.

\textbf{Human label variation} establishes that annotator disagreement ---
particularly in value-laden tasks --- frequently reflects genuine ambiguity or
perspective-dependence rather than annotation error, undermining the assumption
that any single set of reference labels constitutes ground truth~\cite{aroyo2015,
plank2022, uma2021, pavlick2019, basile2021, gordon2022}.

\textbf{Contestability and validity theory} provide both the normative
justification --- contested decisions require procedural legitimacy --- and
the measurement-theoretic foundation --- consequential validity replaces
criterion validity when no stable criterion exists --- for evaluation
approaches that do not depend on reference labels~\cite{messick1995,
alfrink2023contestable, henin2022, hildebrandt2019agonistic, huang2025moderation}.

What is missing is a diagnostic methodology that integrates these insights
operationally: one that takes an assessor and a labelled corpus, determines
whether the reference labels are sufficiently reliable to support
accuracy-based evaluation, and signals when alternative evaluation criteria
are needed.  This dissertation proposes and demonstrates such a methodology,
applying it to four automated AI ethics assessors evaluated on a 40-proposal
corpus with stratified annotations across four professional perspectives.
